\section{RCOUNT Inputs And Report Contract}

The intended inputs are RCOUNT batch manifests, cast vote records, and normalized
summaries. Batch manifests identify the operational grouping: batch, scanner,
drop box, upload, or equivalent source. Cast vote records and summaries provide
the contest totals or ballot-level records from which vote shares and residuals
are derived. Source package hashes bind the report to the exact RCOUNT package
used to compute it.

W.03 instantiates the W-series report contract defined in W.01. The per-record
fields are:

\begin{center}
\small
\begin{tabularx}{\textwidth}{lX}
\toprule
Field & W.03 value or meaning \\
\midrule
\texttt{method\_id} & Registered batch/scanner-effect analytic. \\
\texttt{feature\_set} & Variables used for vote share, residuals, grouping, and baseline adjustment. \\
\texttt{baseline\_population} & Comparable contests, units, modes, times, or groups used for calibration. \\
\texttt{model\_id} & \texttt{batch-effect}. \\
\texttt{unit\_id} & Batch, scanner, drop box, upload, or other operational group. \\
\texttt{score} & Group-level anomaly score or standardized effect estimate. \\
\texttt{p\_value\_or\_rank} & Optional calibration or rank within the baseline population. \\
\texttt{investigation\_note} & Caveat explaining the grouping, comparison, and plausible benign boundary. \\
\bottomrule
\end{tabularx}
\end{center}

The report is about the group, not individual ballots. A scanner effect should
be attributed to the scanner only when the inputs support that grouping. A mail
batch effect should not be described as a precinct effect merely because the
batch contains precinct-coded records. The central discipline is attribution:
identify the level at which the deviation is measured, bind it to source package
hashes, and keep the result in the W-series investigative transcript rather than
V-series certification evidence.
